\section{Introduction}

Thank you for choosing to use \LaTeX\ File Template. If you find any bugs or want some feature not currently included, then my contact information can be found on my GitHub: https://github.com/MatthewEmer.

\separate{7}This document is intended for academics (with special focus on those in the fields of computer science or mathematics) who want to write papers or other related documents in a consistent format which ties in with standard practises from other published papers. This guide consists of setup instructions to download the files and create your first document, before moving onto the custom commands/environments designed to speed up academic writing. Before using this guide, I would suggest that you have some practise with some basic \LaTeX\ commands using a guide such as https://www.overleaf.com/learn.

\separate{7}This guide is accurate as of Version \templateversion.

\section{Document Setup}

\subsection{Project Files}

The actual \LaTeX\ file template is contained in the folder \textit{The Template} and consists of four files: \textit{template.tex}, \textit{\_documentContent.tex}, \textit{\_appendix.tex}, and \textit{template.pdf}. To use the template, all four files must be downloaded and placed in the same folder. 

\subsubsection{\_appendix.tex}

This file allows you to add an optional appendix to the end of your document which is not counted towards your page count. By default it contains a unnumbered part title which has been inserted into the table of contents.  

\subsubsection{\_documentContent.tex}

This file allows you to add all the main content of the document you are trying to write. It will be inserted after the table of contents, but before the appendix (if you choose to add one). By default it contains an example section and definition. Only pages in this file are counted towards the page count of the document.

\subsubsection{template.pdf}

This file contains the output of the \LaTeX\ compiler. If you want to rename this document while still working on it, then delete this version, alongside any build files your compiler has created and change the name of \textit{template.tex}. 

\subsubsection{template.tex}

This file sets up the structure of the document, as well as acting as a compiler for the other two \textit{.tex} files, ensuring they get inserted into the right sections of the document. When creating your document, you need to edit lines 12-15 of this file to create the title page, alongside the document header and footer, as described in \elemRef{tab:Lines to edit.}{Table}.

\inserttable{c|c|c}{Lines to edit.}{\textbf{Line} & \textbf{Variable Name} & \textbf{Description}}
{
    12 & thetitle & The document's title. This will be placed in the header. \\
    13 & thesubtitle & The document's (optional) subtitle. This will be placed in the header. \\
    14 & theauthor & The list of authors. This will be placed in the footer. \\
    15 & thedate & Replace the code if you want a static date rather than date of last compile.
}

\newpage\subsection{Optional Files}

\subsubsection{Storing Images}

If you choose to include figures in your project, then you will need to create a subfolder labelled \textit{Images} where you can place the image files for the project.

\subsubsection{GitHub Repositories}

If you are placing the document within a GitHub repository, you may want to copy over the \textit{.gitignore} from this repository as it is set up to ignore any \LaTeX\ build files.

\section{Custom Commands}

\subsection{General Tools}

\subsubsection{Element Spacing}

\numberedentry{cmd:separate}{Command}{\(\backslash separate\)

\separate{3}The separate command inserts vertical spacing between two elements whilst simultaneously removing the new paragraph indent from the following element. You can specify the amount of spacing as described in \elemRef{tab:Parameters for separate.}{Table} and as seen in \preEnvRef{ex:cmd:separate}{Example}{entrynumber}.

\inserttable{c|c|c|c}{Parameters for separate.}
{\textbf{No.} & \textbf{Name} & \textbf{Description} & \textbf{Example Value}}
{
    1 & Spacing & The number of points of vertical spacing. & 5
}

\noindent\example{cmd:separate}{\(\backslash separate\{3\}\) produces the separation as seen between the Command line and the description of separate.}}

\subsubsection{Text Formatting}

\numberedentry{cmd:colouredText}{Command}{\(\backslash colouredText\)

\separate{3}The colouredText command takes the given text and outputs it in a different colour, as chosen by the writer. This can be seen in \preEnvRef{ex:cmd:colouredText}{Example}{examplenumber} and as described in \elemRef{tab:Parameters for colouredText.}{Table}.

\inserttable{c|c|c|c}{Parameters for colouredText.}
{\textbf{No.} & \textbf{Name} & \textbf{Description} & \textbf{Example Value}}
{
    1 & Colour & The colour to change the text to. & blue \\
    2 & Text & The text you want outputted in colour. & Hello World
}

\noindent\example{cmd:colouredText}{\(\backslash colouredText\{blue\}\{\text{Hello World}\}\) produces \colouredText{blue}{Hello World}.}}

\separate{15}\numberedentry{cmd:comment}{Command}{\(\backslash comment\)

\separate{3}The comment command inserts notes into the document in red so they can be found more easily. This can be seen in \preEnvRef{ex:cmd:comment}{Example}{examplenumber} and as described in \elemRef{tab:Parameters for comment.}{Table}.

\inserttable{c|c|c|c}{Parameters for comment.}
{\textbf{No.} & \textbf{Name} & \textbf{Description} & \textbf{Example Value}}
{
    1 & Text & The text you want outputted in colour. & Hello World
}

\noindent\example{cmd:comment}{\(\backslash comment\{\text{Hello World}\}\) produces \comment{Hello World}.}}

\subsubsection{Referencing Tables and Figures}

\numberedentry{cmd:elemRef}{Command}{\(\backslash elemRef\)

\separate{3}The elemRef command allows you to easily reference tables and figures with clickable links. This is described in \elemRef{tab:Parameters for elemRef.}{Table} and as seen in \preEnvRef{ex:cmd:elemRef}{Example}{examplenumber}.

\inserttable{c|c|c|c}{Parameters for elemRef.}
{\textbf{No.} & \textbf{Name} & \textbf{Description} & \textbf{Example Value}}
{
    1 & Label & The label of the element being linked to. & tab:Parameters for elemRef. \\
    2 & Type & The type of element being linked to. & Table
}

\noindent\example{cmd:elemRef}{\(\backslash elemRef\{\text{tab:Parameters for elemRef.}\}\{\text{Table}\}\) produces the reference \elemRef{tab:Parameters for elemRef.}{Table}.}}

\subsection{Academic Environments}

\subsubsection{Referencing Environments}

\numberedentry{cmd:preEnvRef}{Command}{\(\backslash preEnvRef\)

\separate{3}The preEnvRef command allows you to easily reference academic environments with clickable links before they are defined in the document. This is described in \elemRef{tab:Parameters for preEnvRef and postEnvRef.}{Table} and as seen in \preEnvRef{ex:cmd:preEnvRef}{Example}{examplenumber}.

\inserttable{c|c|c|c}{Parameters for preEnvRef and postEnvRef.}
{\textbf{No.} & \textbf{Name} & \textbf{Description} & \textbf{Example Value}}
{
    1 & Label & The label of the environment being linked to. & ex:cmd:postEnvRef \\
    2 & Type & The type of environment being linked to. & Example \\
    3 & Counter & The name of the counter for that environment. & examplenumber
}

\noindent\example{cmd:preEnvRef}{\(\backslash preEnvRef\{\text{ex:cmd:preEnvRef}\}\{\text{Example}\}\{\text{examplenumber}\}\) produces the reference \preEnvRef{ex:cmd:preEnvRef}{Example}{examplenumber}.}}

\separate{15}\numberedentry{cmd:postEnvRef}{Command}{\(\backslash postEnvRef\)

\separate{3}The postEnvRef command allows you to easily reference academic environments with clickable links after they are defined in the document. This is described in \elemRef{tab:Parameters for preEnvRef and postEnvRef.}{Table} and as seen in \preEnvRef{ex:cmd:postEnvRef}{Example}{examplenumber}.

\separate{5}\example{cmd:postEnvRef}{\(\backslash postEnvRef\{\text{ex:cmd:preEnvRef}\}\{\text{Example}\}\{\text{examplenumber}\}\) produces the reference \postEnvRef{ex:cmd:preEnvRef}{Example}{examplenumber}.}}

\subsection{Mathematics}

\subsection{Computer Science}